\subsection{Unrestricted classification with at most two dimensions}
\label{sec:unbounded-two-dimensions}
The multiplicity bound can be removed entirely when the basis dimensions
assume at most two distinct values. This statement concerns arbitrary
complex categorification, and hence applies simultaneously to all five
existence questions.

\begin{theorem}\label{thm:two-dimensions-unbounded}
Let $R$ be a rank-four fusion ring whose simple Frobenius--Perron dimensions
have at most two distinct values. Then $R$ admits a complex categorification
if and only if it is one of the eleven rings in the following table. Every
listed ring admits a unitary categorification.
\[
\begin{array}{c|l|c}
\text{Number of invertibles}&\text{Fusion ring or defining parameters}
 &\text{Number of rings}\\\hline
4&\mathbb Z[C_4],\ \mathbb Z[C_2\times C_2]&2\\
3&\mathrm{NG}(C_3,m),\quad m\in\{0,2,3,6\}&4\\
2&gX=Y,\ X=X^*,\ X^2=1+aX+bY,\\[-3pt]
 &\hspace{1em}(a,b)\in\{(1,0),(1,1),(2,2)\}&3\\
2&S(1)&1\\
2&K_0(\operatorname{Rep}(D_{10}))&1
\end{array}
\]
Here $D_{10}$ is the dihedral group of order ten, $g$ is the nontrivial
invertible in the third row, and $\mathrm{NG}(C_3,m)$ denotes the near-group
ring with the three group elements and one object $X$ satisfying
$X^2=\sum_{g\in C_3}g+mX$. The third-row rings are unchanged up to based-ring
isomorphism when $a$ and $b$ are interchanged.
\end{theorem}
\begin{proof}
The pointed case is precisely the first row. Suppose henceforth that the
noninvertible dimensions all equal $d>1$, and let $G$ be the group of
invertible basis elements. Schopieray's theorem
\cite[Lemma~3.1]{Schopieray} reduces general categorifiability to the
existence of a pseudo-unitary categorification.

We first show that $|G|\geq2$. Otherwise, choose noninvertible basis
elements $X,Y$ with $Y\ne X^*$. Their product has no unit summand, so
$d^2=kd$ for a nonnegative integer $k$ and $d=k$ is an integer. On the
other hand, $XX^*$ has exactly one unit summand and no other invertible
summand, giving $d^2=1+\ell d$ for an integer $\ell$. This forces $d\mid1$,
a contradiction.

If $|G|=3$, its action fixes the unique noninvertible object. Frobenius
reciprocity forces the near-group fusion rule in the table. Larson's
classification \cite[Theorem~1.1]{Larson} gives precisely
$m\in\{0,2,3,6\}$ in the pseudo-unitary setting. The unitary realizations
are recorded in Table~\ref{B-tab:models} of \cite{Finite}.

It remains to consider $G=\{1,g\}$ and the two noninvertible objects $X,Y$.
If $g$ interchanges them, the duality either fixes both or interchanges
them. In the self-dual case, Frobenius reciprocity gives
\[
 X^2=Y^2=1+aX+bY,\qquad XY=g+bX+aY.
\]
The pseudo-unitary restriction
\cite[Theorem~2.4]{EIP} gives, up to interchanging $X$ and $Y$, the
three nonpointed parameter pairs in the table. In the non-self-dual case,
$X^2=Y^2$ is fixed by duality, so its coefficients on $X,Y$ agree. Thus
$R=S(m)$ for some $m\geq1$. \cite[Theorem~\ref{B-thm:nonselfdual-quadratic}]{Finite}
leaves exactly $m=1$. All four rings have unitary realizations as stated
in the census table.

Finally suppose $g$ fixes both noninvertibles. For $X\ne Y$, the product
$XY^*$ has no invertible summand: the multiplicity of $h\in G$ equals
$[hY:X]=0$. Hence $d^2=kd$, so $d$ is an integer. The product $XX^*$
contains both $1,g$ once, giving $d^2=2+\ell d$. Therefore $d=2$.
If $X^*=Y$, then $XY$ is invariant under duality, so its noninvertible
coefficients agree. Its dimension equation would be $4=2+4c$, impossible
for integral $c$. Thus both objects are self-dual. Write
\[
 X^2=1+g+aX+bY,\quad XY=bX+cY,\quad
 Y^2=1+g+cX+eY.
\]
The dimension equations become
$a+b=1$, $b+c=2$, and $c+e=1$. Nonnegativity forces
$(a,b,c,e)=(0,1,1,0)$, which is the representation ring of $D_{10}$.
This exhausts the cases. The known unitary realizations prove sufficiency.
\end{proof}

\begin{remark}[What this leaves open]
The theorem is independent of multiplicity, but its dimension hypothesis
is essential to the argument. In particular, it does not settle arbitrary
self-dual rank-four rings with three or four distinct simple dimensions.
The unrestricted exclusion of the family $R_n$ for all $n\geq2$ supplies
one further complete infinite family; it does not prove that these families
exhaust the remaining rings.
\end{remark}
